% ===== tr-07 块3：第 5 章 pp.8--10（原文页 8--10）=====

% [原文 p.8]
近来人们认识到，在复维数三的情形存在几乎典范的 Ooguri--Vafa 型度量构造，它们具有所预测的正顶点与负顶点的拓扑与复结构，并可由（非线性的）广义 Gibbons--Hawking ansatz 构造出来 [54]。这些 Ooguri--Vafa 型度量的渐近几何在一般区域中与半平坦度量相配。正顶点度量包含一个以 $\C^3$ 上广义 Taub--NUT 型度量为模型的局部区域，这类似于 Ooguri--Vafa 度量包含一个以 Taub--NUT 度量为模型的区域。

下列问题广为公开：

\begin{question}{5.1}
这类 Ooguri--Vafa 型度量是否会作为任何紧 Calabi--Yau 流形在大复结构极限附近的爆破极限出现？
\end{question}

\begin{question}{5.2}
能否对三维流形在大复结构附近给出其 CY 度量的粘合描述，例如在五次超曲面的情形？
\end{question}

\begin{question}{5.3}
在这类 Ooguri--Vafa 型度量的 $C^\infty$-小扰动上，可以出现哪些种类的特殊拉格朗日？
\end{question}

% [Tutor Note: 原文这三问编号为 Question 1/2/3（与 p.9 的 Question 4 连成一列），按本项目的定理类编号规则写成 5.1--5.3；这些编号与本章公式编号 5.1--5.6 属于不同的计数体系，正文中引用时需靠“问题/式”区分。]

\subsection{强 SYZ 猜想与弱 SYZ 猜想}

由于 SYZ 猜想的分析困难，尤其是具有大黎曼曲率的非一般区域带来的困难，文献中发展出了许多种对猜想的解读，其目标也各有分歧。

\begin{itemize}
\item （软版本）为了应用于同调镜像对称，人们主要感兴趣的是构造拉格朗日纤维化而不要求 $\mathrm{Im}(\Omega)|_L=0$，并由此构造镜像流形 $X^\vee$、证明范畴层面的预言 $D^bCoh(X^\vee)\simeq D^\pi Fuk(X,\omega)$。这种观点把 Calabi--Yau 流形的辛数据与全纯数据分离开来，具有拓扑/代数的味道。

\begin{remark}{5.1}
通常在同调镜像对称的讨论中被略去的特殊拉格朗日条件，据信与 Bridgeland 稳定性条件有关，后者是 D-膜上 BPS 条件的一种流行的范畴式解释。
\end{remark}

\item （强度量版本）在紧 Calabi--Yau 流形大复结构极限附近，找到一个整体的特殊拉格朗日环面纤维化。
\item （弱度量版本）对某一类合适的紧 Calabi--Yau 流形证明：在大复结构极限附近，特殊拉格朗日 $T^n$-纤维化存在于一个占 Calabi--Yau 测度至少 $99\%$ 的大子集上。更精确地说，该百分比在极限下收敛到 $100\%$。
\end{itemize}

% [Tutor Note: 弱版本中的“$99\%$”与“收敛到 $100\%$”是同一件事的两种说法：前者是对某个足够接近极限的成员给出的定量下界，后者是沿族的极限陈述；原文把两者并列，不能理解为“99% 之后另有 100% 的结论”。]

% [原文 p.9]
度量版本对 Calabi--Yau 度量高度敏感，而后者与软版本相比是刚性得多的结构。它们符合 SYZ 原文的偏微分方程精神，而软版本则更接近 SYZ 猜想的镜像对称动机。

强度量版本也许最忠实于 SYZ 原文的初衷。它的主要证据来自超 Kähler 度量这一特殊情形，已在第 5.2.3 节中提到。其他边缘性的证据来自许多例子中拉格朗日纤维化的构造，那里忽略了 $\mathrm{Im}(\Omega)|_L=0$ 这一条件 [58]。然而有许多微妙之处困扰着这一强版本（已在第 5.2.4 节中提到），使该猜想变得非常艰巨，相比之下支持它的证据显得不足。值得注意的是，特殊拉格朗日奇点尚未被充分理解，而且我们不知道有任何论证能明确排除特殊拉格朗日在具有大黎曼曲率的非一般区域中彼此相交。作为供今后思考的素材，一个绕过这些可能反对意见的稍弱版本是

\begin{question}{5.4}
给定紧 Calabi--Yau 流形 $X$，当它足够接近大复结构极限时，是否存在一族 $n$ 参数的特殊拉格朗日流，其支撑扫过 $X$ 上的所有点？
\end{question}

另一方面，弱度量版本只涉及一般区域，因而比强版本更易处理。它符合更谨慎的期望，即特殊拉格朗日纤维化只是一种极限现象。本综述的主体将聚焦于这一弱版本。

\section{大复结构极限}

SYZ 原文并未精确给出大复结构极限的概念，而当前文献中有若干互不等价的解释。我们将把大复结构极限放在极化退化的框架中。直观地说，极化退化是指：把辛结构固定在一个整类中，并让复结构在一个代数单参数族中变化，使其在极限中变成奇异的。大复结构极限则是附加的要求------退化“尽可能剧烈”。

我们在 $\C$ 上工作。为了铺陈背景，先作如下约定：
\begin{itemize}
\item 设 $S$ 是一条光滑仿射代数曲线，带有一个点 $0\in S$。一个代数退化族由一个淹没的射影态射 $\pi:X\to S\setminus\{0\}$ 给出，其纤维 $X_t$（$t\in S\setminus\{0\}$）是光滑连通、$n$ 维的。这与在穿孔形式圆盘 $\operatorname{Spec}(K)$（$K=\C(\!(t)\!)$）上的形式框架形成对照。代数退化经基变换诱导出形式退化。
\end{itemize}

% [原文 p.10]
\begin{itemize}
\item 一个极化由 $X$ 上的一个丰沛线丛 $L$ 给出。这确定了 Kähler 类，至多相差重新标度的约定。
\item 若典范丛 $K_X$ 有一个平凡化截面 $\Omega$，则称 $\pi$ 是 Calabi--Yau 流形的退化族。在 $0\in S$ 附近的一个小圆盘 $\mathbb{D}_t$ 上，$\Omega$ 通过 $\Omega=dt\wedge\Omega_t$ 诱导出 $X_t$ 上的全纯体积形式 $\Omega_t$。$X_t$ 上的归一化 Calabi--Yau 测度是概率测度
\begin{equation}
d\mu_t = \frac{\Omega_t\wedge\overline{\Omega_t}}{\int_{X_t}\Omega_t\wedge\overline{\Omega_t}}. \tag{5.5}
\end{equation}
$X_t$ 上的 Calabi--Yau 度量 $\omega_{CY,t}$ 是类 $\frac{1}{|\log|t||}c_1(L)$ 中唯一的 Kähler 度量，使得
\begin{equation}
\frac{\omega_{CY,t}^n}{\int_{X_t}\omega_{CY,t}^n} = d\mu_t. \tag{5.6}
\end{equation}
\item 称 $\pi:X\to S\setminus\{0\}$ 是 Calabi--Yau 流形的大复结构极限，若本质骨架具有最大维数 $n$（下文解释），且该退化族在 $S$ 上容许一个半稳定 snc 模型。
\end{itemize}

% [Tutor Note: 式 (5.6) 中分母是 $\int_{X_t}\omega_{CY,t}^n$（体积的 $n$ 次幂积分），分子与分母同除后左端是关于 $\omega_{CY,t}$ 的概率测度；式 (5.5) 与 (5.6) 联立即“Calabi--Yau 度量按归一化体积测度等于 $\Omega_t\wedge\overline{\Omega_t}$ 所定测度”，这正是 Calabi--Yau 条件的度量表述。抄写时易把 (5.6) 左端误作 $\omega_{CY,t}^n/\int\Omega_t\wedge\overline{\Omega_t}$。]

% [Tutor Note: 类 $\frac{1}{|\log|t||}c_1(L)$ 中的绝对值有两层：内层 $\log|t|$ 取 $t$ 的模，外层再取绝对值；$t\to0$ 时 $\log|t|\to-\infty$，故该 Kähler 类按 $1/|\log|t||$ 的速度趋于零------这是度量坍缩的标度。]

\begin{remark}{5.2}
在 $0\in S$ 处填上中心纤维，将涉及为 $\pi:X\to S$ 选取一个模型，即一个正规、平坦、射影的 $S$-概形 $\mathcal{X}$，连同在穿孔曲线 $S\setminus\{0\}$ 上与 $X$ 的一个同构。若 $\mathcal{X}$ 光滑，且 $0\in S$ 上的中心纤维是 $\mathcal{X}$ 中的简单正规交叉除子，使得各除子之交为不可约的或为空，则称之为 snc 模型。若进一步中心纤维是约化的，则称之为半稳定 snc 模型。模型也可以在形式圆盘上类似地定义。snc 模型的存在性是 Hironaka 消解定理的推论。它们高度不唯一。由半稳定约化定理 [44, 第 2 章]，经过到另一条光滑代数曲线 $S'$ 的有限基变换后，我们总能对退化族 $X\times_S(S'\setminus\{0\})$ 找到某个半稳定 snc 模型，因此半稳定 snc 模型的存在性并非实质性假设。这里的一切都是拟射影的。模型的选取非常有用，但对退化的 CY 度量而言并非内蕴的。
\end{remark}

\begin{example}{5.3.1}
大复结构极限的一个典型例子是 Calabi--Yau 超曲面的 Fermat 族
\[
X_t = \Big\{Z_0Z_1\dots Z_{n+1} + t\sum_{0}^{n+1}Z_i^{n+2}=0\Big\}\subset\C\PP^{n+1}.
\]
当 $t\to0$ 时，代数极限是 $n+2$ 个射影平面的并。直观上，中心纤维高度可约，退化非常剧烈。
\end{example}

% [Tutor Note: 原文 Example 3.1 说极限是 “$n+2$ 个 projective planes 的并”；对一般的 $n$，$Z_i=0$ 截出的是 $\C\PP^{n+1}$ 中的射影超平面 $\cong\C\PP^n$（仅当 $n=2$ 时才是“平面”）。译文照原文作“射影平面”，读者应理解为 $\C\PP^n$。$n+2$ 个因子与 $Z_0,\dots,Z_{n+1}$ 共 $n+2$ 个变量相符。]

\begin{example}{5.3.2}
考虑一族退化为结点 K3 曲面的四次 K3 曲面。这是一个典型的极化退化，但它不是大复结构极限。
\end{example}
